\section*{अध्याय \ref{ch:b2:plant-flowering} --- जननीय विकास और पुष्पन का नियंत्रण}

\begin{solution}{exo:b2:plant-flowering:1}
कायिक \omterm{def:b2:plant-meristems:meristem}{विभज्योतक}: कक्षस्थ कलिकाओं सहित पत्तियाँ बनाता है, अनिश्चित काल तक
(अनिश्चित)। \omterm{def:b2:plant-flowering:transition}{पुष्पक्रम विभज्योतक}: सहपत्र और उनके कक्षों में पुष्पीय
\omterm{def:b2:plant-meristems:meristem}{विभज्योतक} बनाता है, और अनेक जातियों में अनिश्चित रहता है। पुष्पीय
\omterm{def:b2:plant-meristems:meristem}{विभज्योतक}: चारों चक्र बनाकर रुक जाता है --- यानी निश्चित।
\end{solution}

\begin{solution}{exo:b2:plant-flowering:2}
अल्पदिवसी पौधे तब फूलते हैं जब दिन किसी क्रांतिक लंबाई से छोटा हो
(गुलदाउदी, सोयाबीन, धान); दीर्घदिवसी पौधे तब जब वह
लंबा हो (पालक, गेहूँ, \emph{Arabidopsis}); और दिन-उदासीन पौधे
कोई आकार पा लेने पर फूलते हैं (टमाटर, मक्का)। कोई \omterm{prop:b2:plant-flowering:photoperiod}{अल्पदिवसी पौधा}
अटूट रात की लंबाई नापता है।
\end{solution}

\begin{solution}{exo:b2:plant-flowering:3}
वन्य प्ररूप: बाह्यदल, दलपत्र, \omterm{def:b2:angiosperm-reproduction:flower}{पुंकेसर}, \omterm{def:b2:angiosperm-reproduction:flower}{अंडप}। A नहीं: \omterm{def:b2:angiosperm-reproduction:flower}{अंडप}, \omterm{def:b2:angiosperm-reproduction:flower}{पुंकेसर},
\omterm{def:b2:angiosperm-reproduction:flower}{पुंकेसर}, \omterm{def:b2:angiosperm-reproduction:flower}{अंडप}। B नहीं: बाह्यदल, बाह्यदल, \omterm{def:b2:angiosperm-reproduction:flower}{अंडप}, \omterm{def:b2:angiosperm-reproduction:flower}{अंडप}। C नहीं: बाह्यदल,
दलपत्र, दलपत्र, बाह्यदल, और वह \omterm{def:b2:angiosperm-reproduction:flower}{पुष्प} भीतर दोहरा जाता है।
\end{solution}

\begin{solution}{exo:b2:plant-flowering:4}
सप्ताहों की ठंड और नमी (शीत ऋतु बीत चुकी); प्रकाश (वह \omterm{def:b2:angiosperm-reproduction:seed}{बीज} सतह के
पास पड़ा है); बदलते ताप (नंगी मिट्टी, कोई छत्र नहीं);
वर्षा से बहाव (वर्षा आ चुकी); खुरचन (किसी आँत से गुज़रना या
मिट्टी में घिसाव); और धुआँ तथा ताप (किसी आग ने ज़मीन साफ़ कर दी है)।
\end{solution}

\begin{solution}{exo:b2:plant-flowering:5}
16/8: 8 घंटे की रात, क्रांतिक 10 से कम: नहीं। 12/12: हाँ।
12 प्रकाश, 6 अंधकार, झलक, 6 अंधकार: 6-6 घंटे की दो रातें: नहीं।
8 प्रकाश, 4 अंधकार, 4 प्रकाश, 8 अंधकार: सबसे लंबी रात 8 घंटे: नहीं।
\end{solution}

\begin{solution}{exo:b2:plant-flowering:6}
\omterm{def:b2:plant-flowering:florigen}{फ़्लोरिजन} ऐसा ग्राफ़्ट-संचारणीय संकेत है जो पत्ती में बनता है और
अल्पदिवसी तथा दीर्घदिवसी जातियों में एक ही है --- भेद इसमें है कि
वह पत्ती उसे कब बनाती है, यह नहीं कि वह क्या बनाती है। फ्लोएम काट दिया
जाए तो वह ग्राफ़्ट पुष्पन प्रेरित नहीं करता: यानी वह संकेत
फ्लोएम में यात्रा करता है।
\end{solution}

\begin{solution}{exo:b2:plant-flowering:7}
$n = \ln(1/0.15)/0.05 = 1.90/0.05 = 38$ ठंडे दिन। उन लहरों का कुल
$47$ है: वह पौधा तीसरी लहर के 11वें दिन वसंतीकृत हो जाता है।
$\alpha = 0.02$ के साथ उसे 95 दिन चाहिए होते और वह नहीं होता।
\end{solution}

\begin{solution}{exo:b2:plant-flowering:8}
उस दमनकारी का जीन ऐसे क्रोमैटिन चिह्नों से चुप किया जाता है जो हर
समसूत्री विभाजन पर नक़ल होते हैं, इसलिए उस प्ररोह की हर कोशिका वह शीत ऋतु
याद रखती है; जबकि जनन वंशक्रम और आरंभिक भ्रूण में वे चिह्न मिटा दिए जाते
हैं और वह जीन फिर अभिव्यक्त होता है। उपयोगी इसलिए कि गर्मी में झड़े \omterm{def:b2:angiosperm-reproduction:seed}{बीज}
को अपने जनक की पुष्पन-अनुमति विरासत में नहीं मिलनी चाहिए: उसे अपनी ही
शीत ऋतु की प्रतीक्षा करनी चाहिए।
\end{solution}

\begin{solution}{exo:b2:plant-flowering:9}
R: अंकुरित (फ़ाइटोक्रोम दूर-लाल-अवशोषी रूप, यानी Pfr, में बदला, जो
सक्रिय है)। R फिर FR: नहीं (वापस Pr में बदला)।
R, FR, R: अंकुरित। R, FR, R, FR: नहीं। केवल अंतिम झलक गिनी जाती है,
क्योंकि हर रूपांतरण पूरा होता है और वह \omterm{def:b2:angiosperm-reproduction:seed}{बीज} उस रंजक की अंतिम
अवस्था पढ़ता है।
\end{solution}

\begin{solution}{exo:b2:plant-flowering:10}
$p = 0.7$: $r = 1.28$, $1.42$, $1.40$, $-\infty$, यानी $g = 0.4, 0.6, 0.8,
1$ के लिए: सर्वोत्तम $g = 0.7$ के पास। $p = 0.95$: $1.99$, $2.33$, $2.55$,
$-\infty$: इष्टतम $g = 1$ की ओर खिसक जाता है (लगभग $0.948$), यानी
नाम भर का भंडार। रेगिस्तान, जहाँ अच्छे वर्ष विरल हैं, भारी
\omterm{def:b2:angiosperm-reproduction:seed}{प्रसुप्ति} और बड़े \omterm{def:b2:angiosperm-reproduction:seed}{बीज} बैंक के पक्ष में हैं; जबकि घास-मैदान, जहाँ
अधिकांश वर्ष अच्छे हैं, लगभग सब कुछ अंकुरित करने के पक्ष में।
\end{solution}

\begin{solution}{exo:b2:plant-flowering:11}
चारों चक्रों में B: दलपत्र, दलपत्र, \omterm{def:b2:angiosperm-reproduction:flower}{पुंकेसर}, \omterm{def:b2:angiosperm-reproduction:flower}{पुंकेसर}। चारों चक्रों में C (C
A को दबाता है): \omterm{def:b2:angiosperm-reproduction:flower}{अंडप}, \omterm{def:b2:angiosperm-reproduction:flower}{पुंकेसर}, \omterm{def:b2:angiosperm-reproduction:flower}{पुंकेसर}, \omterm{def:b2:angiosperm-reproduction:flower}{अंडप}। न A न B: हर जगह अकेला C
--- यानी चारों चक्रों में \omterm{def:b2:angiosperm-reproduction:flower}{अंडप}। C \omterm{def:b2:plant-meristems:meristem}{विभज्योतक} को समाप्त भी करता
है; उसके बिना वह केंद्र अनिश्चित बना रहता है और \omterm{def:b2:angiosperm-reproduction:flower}{पुष्प}
के भीतर \omterm{def:b2:angiosperm-reproduction:flower}{पुष्प} बनाता है।
\end{solution}

\begin{solution}{exo:b2:plant-flowering:12}
दिन-लंबाई रातों के आरपार इस बात से समाकलित होती है कि प्रकाश उस घड़ी से
मेल खाए, और प्रेरण को ऐसी कई रातें चाहिए; ठंड जमा होते क्रोमैटिन चिह्नों
के रूप में समाकलित होती है; और \omterm{def:b2:angiosperm-reproduction:seed}{बीज} बैंक
अंकुरण फैलाकर वर्षों के आरपार समाकलन करता है। कोई समाकलक
किसी अकेले असामान्य दिन, जनवरी के किसी पिघलाव, या एक बुरे वर्ष को
अनदेखा कर देता है; जबकि तात्क्षणिक मान पर कोई देहली उस पौधे को फ़रवरी के
किसी गर्म सप्ताह में फुला देती, या पहली ही वर्षा में हर \omterm{def:b2:angiosperm-reproduction:seed}{बीज} अंकुरित कर देती।
\end{solution}

\begin{solution}{pb:b2:plant-flowering:1}
\textbf{1.} 8 घंटे की रात क्रांतिक 9.5 घंटे से छोटी है:
कोई पुष्पन नहीं। अगस्त में बेचने के लिए उस उगाने वाले को रातें काले
कपड़े से कृत्रिम रूप से लंबी करनी पड़ेंगी।
\textbf{2.} 12 घंटे की रात। 15 अगस्त को खिलने का अर्थ है प्रेरण उससे
8 सप्ताह पहले, यानी लगभग 20 जून, पूरा हो; इसलिए वे बारह प्रेरक
रातें लगभग 8 जून को आरंभ होनी चाहिए।
\textbf{3.} उस फ़सल को हर रात के बीचोंबीच कुछ मिनट रोशनी दीजिए:
14.5 घंटे की रात लगभग 7-7 घंटे की दो रातें बन जाती है, और दोनों
क्रांतिक लंबाई से कम। झलक इसलिए पर्याप्त है कि फ़ाइटोक्रोम
मिनटों में बदल जाता है और वह पौधा केवल अटूट अंधकार की लंबाई नापता है।
\textbf{4.} 21:00 पर दी गई झलक 21:00 से 07:30 तक,
यानी 10.5 घंटे, का अंधकार छोड़ देती है --- जो 9.5 से लंबा है:
इसलिए वह खंड प्रेरण करता है। वह विराम ऐसा पड़ना चाहिए कि कोई भी खंड
क्रांतिक लंबाई से आगे न जाए।
\textbf{5.} एक छूटी हुई रात बारह लगातार प्रेरक रातों की उस दौड़ को तोड़
देती है: गिनती फिर से आरंभ होती है और पुष्पन लगभग उतने ही दिन टल जाता है
जितने जमा हो चुके थे। जबकि किसी कॉकलबर को एक ही प्रेरक रात चाहिए और वह
पहले ही प्रतिबद्ध हो चुका होता।
\textbf{6.} फ़ाइटोक्रोम हरा प्रकाश अवशोषित नहीं करता, इसलिए वह पौधा
उस झलक को देखता ही नहीं; लाल उसे सक्रिय रूप में बदल देता है, जो
दिन की तरह पढ़ा जाता है; और तुरंत बाद का दूर-लाल उसे वापस बदल देता है,
इसलिए वह रात टूटती नहीं।
\textbf{7.} वह \omterm{prop:b2:plant-flowering:photoperiod}{दीर्घदिवसी पौधा} फूल जाता है: \omterm{def:b2:plant-flowering:florigen}{फ़्लोरिजन} उस ग्राफ़्ट
संधि को फ्लोएम में पार कर जाता है। यह दोनों दिन-लंबाई वर्गों में साझा
एक गतिशील, ग्राफ़्ट-संचारणीय संकेत की पहचान कराता है।
\textbf{8.} $n = \ln 10/0.06 = 38.4$: यानी 39 ठंडे दिन।
\textbf{9.} नवंबर में 10, जनवरी के अंत तक 35, और 39वाँ
ठंडा दिन फ़रवरी के चौथे ठंडे दिन पर पड़ता है: यानी फ़रवरी के आरंभ में
वसंतीकरण।
\textbf{10.} छह दिन: $f = e^{-0.36} = 0.70 > 0.10$: वह पूरी गर्मी
कायिक बना रहता है और कोई दाना नहीं देता --- शीत-गेहूँ
शरद में ही बोया जाना चाहिए।
\textbf{11.} उस दमनकारी के बिना वह बिना ठंड के फूल जाता है: उसे
वसंत में बोकर उसी ग्रीष्म काटा जा सकता है --- यानी एक वसंत गेहूँ।
\textbf{12.} उस दमनकारी के क्रोमैटिन पर चुप्पी के चिह्न हर समसूत्री
विभाजन पर नक़ल होते हैं, इसलिए वह शीर्ष वसंत भर याद रखता है; जबकि
\omterm{def:b2:meiosis-heredity:meiosis}{अर्धसूत्री विभाजन} और भ्रूण में वे चिह्न मिट जाते हैं, और हर पीढ़ी को
अपनी शीत ऋतु स्वयं गिननी पड़ती है।
\textbf{13.} वह स्मृति उन्हीं कोशिकाओं में होनी चाहिए जो
\omterm{def:b2:angiosperm-reproduction:flower}{पुष्प} बनाएँगी और जो शीत ऋतु भर बनी रहती हैं --- यानी \omterm{def:b2:plant-meristems:meristem}{विभज्योतक} में --- और
वह उनके क्रोमैटिन की एक अवस्था है, कोई गतिशील संकेत नहीं; जबकि जो
पत्तियाँ दिन-लंबाई नापती हैं वे झड़ जाती हैं, और एक प्रोटीन भेजती हैं।
\textbf{14.} वर्ग C: बाह्यदल, दलपत्र, दलपत्र, बाह्यदल (A चौथे चक्र में
फैल जाता है और वह \omterm{def:b2:angiosperm-reproduction:flower}{पुष्प} भीतर दोहरा जाता है)।
\textbf{15.} कोई \omterm{def:b2:angiosperm-reproduction:flower}{पुंकेसर} नहीं, इसलिए कोई \omterm{def:b2:plant-life-cycles:heterospory}{पराग} नहीं; पर जहाँ C
कमज़ोर रूप से सक्रिय है वहाँ कुछ \omterm{def:b2:angiosperm-reproduction:flower}{अंडप} बन जाते हैं, इसलिए कभी-कभार \omterm{def:b2:angiosperm-reproduction:seed}{बीज} बँध जाते हैं।
\textbf{16.} वर्ग A: \omterm{def:b2:angiosperm-reproduction:flower}{अंडप}, \omterm{def:b2:angiosperm-reproduction:flower}{पुंकेसर}, \omterm{def:b2:angiosperm-reproduction:flower}{पुंकेसर}, \omterm{def:b2:angiosperm-reproduction:flower}{अंडप}।
\textbf{17.} C A को दबाता है: चक्र 1 \omterm{def:b2:angiosperm-reproduction:flower}{अंडप} बन जाता है, चक्र 2 (B और C)
\omterm{def:b2:angiosperm-reproduction:flower}{पुंकेसर} --- यानी \omterm{def:b2:angiosperm-reproduction:flower}{अंडप}, \omterm{def:b2:angiosperm-reproduction:flower}{पुंकेसर}, \omterm{def:b2:angiosperm-reproduction:flower}{पुंकेसर}, \omterm{def:b2:angiosperm-reproduction:flower}{अंडप}, वही A-उत्परिवर्ती \omterm{def:b2:angiosperm-reproduction:flower}{पुष्प}।
\textbf{18.} ABC जीन \omterm{prop:b2:cell-differentiation:transcriptional}{अनुलेखन कारकों} का एक ही कुल हैं
जो सारे सपुष्पक पौधों के साझा पूर्वज से विरासत में मिला है; वह
मॉड्यूल गुलाब, स्नैपड्रैगन और \emph{Arabidopsis} के अलग होने से पहले ही
बैठ चुका था, और हर एक ने उसे बनाए रखा है।
\textbf{19.} A या C के बिना B कोई पुष्पीय अंग निर्दिष्ट नहीं करता:
यानी कोई पत्ती जैसा या सहपत्र जैसा अंग।
\textbf{20.} $p = 0.8$: $r = 1.44$ ($g = 0.5$), $1.59$ ($0.7$), $1.55$
($0.9$)।
\textbf{21.} इन तीनों में $g = 0.7$ (असल इष्टतम
$0.8$ के पास पड़ता है); और $g = 1$ पर बुरे वर्ष का गुणक शून्य है तथा वह
वंशक्रम पहले ही बुरे वर्ष में मिट जाता है।
\textbf{22.} $p = 0.5$: $0.51$, $0.41$, $-0.03$: $g = 0.5$ सबसे अच्छा
है, और इससे भी कम और अच्छा होता --- यानी रेगिस्तान में अधिक \omterm{def:b2:angiosperm-reproduction:seed}{प्रसुप्ति}।
\textbf{23.} उस उगाने वाले के \omterm{def:b2:angiosperm-reproduction:seed}{बीज} गाड़े नहीं बल्कि सँजोए जाते हैं और
निश्चित रूप से बच जाते हैं; उसे खेत के भीतर दाँव बाँटने की ज़रूरत नहीं,
पर किसी विफल वर्ष के बाद फिर बोने को खलिहान में एक भंडार रखना ही पड़ेगा
--- यानी वही भंडार, जो मिट्टी के बजाय उसके पास रहता है।
\textbf{24.} कोई छोटा \omterm{def:b2:angiosperm-reproduction:seed}{बीज} एक-दो सेंटीमीटर के प्ररोह लायक़ भंडार ढोता
है; मिट्टी के नीचे अँधेरे में अंकुरित होकर वह प्रकाश तक पहुँचने से पहले
ही मर जाता, इसलिए वह प्रकाश की प्रतीक्षा करता है --- और वह उस तक
केवल सतह के पास ही पहुँचता है। जुताई गड़े \omterm{def:b2:angiosperm-reproduction:seed}{बीजों} को ऊपर ले आती है और
खरपतवार की एक बाढ़ ला देती है।
\textbf{25.} क्रांतिक रात 9.5 घंटे, और 15 अगस्त को खिलने के लिए लगभग
8 जून से ढकना; 39 ठंडे दिन; और सर्वोत्तम $g \approx 0.8$ जब $p = 0.8$
हो तथा लगभग $0.5$ जब $p = 0.5$ हो।
\end{solution}
